\begingroup
\clearpage\SourceLeaf{330b}\ArithmeticHead{323}{R}{{\scriptsize BEMERKUNGEN ZUR THEORIE DER TRANSZENDENTEN FUNKTIONEN.}}
\begin{center}\bfseries\textls[120]{Erläuterungen zu [III.], S. 311---320.}\end{center}
\par\indent Diese Abhandlung besteht aus zwei zu verschiedenen Zeiten entstandenen Teilen. \textls[120]{Der erste}, die artt. [1.]---[4.] und noch einige beim Abdruck weggelassene Formeln enthaltende, steht auf den S. 172---177 des Handbuchs 19, Be; er stammt wohl aus dem Jahre 1834, da auf S. 184 des Handbuchs (gedruckt Werke V, S. 609, art. [7.]) das Datum »1835, Januar 23« steht. Im art. [1.] ist die Aufgabe berührt, für eine geschlossene Linie $L$, deren Punkte die Koordinaten $x,y$ haben, $x+yi$ als Funktion von $\cos\varphi+i\sin\varphi$ darzustellen, d. h. diese Linie auf den Kreis mit dem Halbmesser Eins abzubilden. In den artt. [3.], [4.] sucht \textsc{Gauss} für den Fall, wo $L$ eine Ellipse ist, diese Aufgabe in der Weise zu lösen, daß er für $T=x+yi$ die Reihe [1] mit reellen Koefficienten ansetzt und die Koefficienten so bestimmt, daß die Reihe [1] und ihr konjugierter Wert [2] in die Gleichung [3], die nichts anderes ist als
\[\frac{x^2}{a^2}+\frac{y^2}{b^2}=1,\]
\noindent eingesetzt, diese für $t=\cos\varphi+i\sin\varphi$ befriedigen. Nach Einführung von $\lambda$ statt $k$ durch die Gleichung [4] ergibt sich die Entwicklung [5], von der ausgehend dann durch Vermittlung der Gleichung [5a] der Zusammenhang zwischen $t$ und der durch [6] definierten Größe $u$ hergestellt wird. Wie \textsc{Gauss} dazu gelangt ist, den Faktor***) von $u+\lambda u^{-1}$ auf der rechten Seite der Gleichung [5a] gerade so zu wählen, geht aus der Aufzeichnung nicht hervor; wahrscheinlich hat ihn eine numerische Induktion dabei geleitet. Den Zusammenhang dieser Reihe, sowie der Reihe [7] mit der Theorie der elliptischen Funktionen findet \textsc{Gauss} erst nach etwa fünf Jahren (siehe die Bemerkung am Schluß des art. [4.]). Zu der durch die Gleichung [5a] gegebenen Darstellung von $x+yi$, wenn $x,y$ die Koordinaten eines Ellipsenpunktes sind, vergleiche man den art. [2.] und die, Werke VIII, S. 341 abgedruckte --- wohl aus späterer Zeit stammende --- Aufzeichnung. Mit dieser Form der Darstellung einer Ellipse beginnt auch der \textls[120]{zweite Teil} unserer Abhandlung im art. [5.].
\par\bigskip\noindent\rule{.25\linewidth}{.4pt}\par\smallskip
{\small\par\indent ***) \textsc{Gauss} hat später ein $\ast$ hingesetzt, um diese Reihe anführen zu können, siehe den art. [7.] S. 318.\par}
\par\medskip\hfill41\textsuperscript{*}\par
\clearpage\SourceLeaf{331}\ArithmeticHead{324}{L}{ANALYSIS.\quad NACHLASS.}
\par\indent Dieser zweite Teil ist auf S. 222---226 desselben Handbuchs 19, Be aufgezeichnet, seine Abfassungszeit ist in dem Hinweis am Schluß des art. [4] genau angegeben. Für eine Ellipse, deren Halbachsen zu $1+\lambda$ und $1-\lambda$ angenommen werden, wird gefragt, mit welcher Dichte die Masse $2\pi A$ auf der Peripherie ausgebreitet werden muß, damit ihr \textls[120]{Potential}*) in den Punkten der Ellipse übereinstimmt mit dem Potential der im Ellipsenmittelpunkt aufgespeicherten Masse. Die Differenz dieser beiden Potentiale ist also nach heutiger Ausdrucksweise nichts anderes, als die zu dem Mittelpunkt als Pol gehörige \textsc{Green}sche Funktion der Ellipsenfläche, multipliziert mit $2\pi A$. Im art. [6.] bildet \textsc{Gauss} die Differenz der \textls[120]{vollständigen Potentiale}, d. h. er denkt sich jedes der beiden Potentiale \linebreak durch Hinzufügen der sogenannten konjugierten Funktion zu einer monogenen Funktion der komplexen Veränderlichen $x+yi$ ergänzt und diese beiden Funktionen von einander subtrahiert; setzt man diese Differenz gleich $\log t$, so liefert $t$ offenbar die in den kleinsten Teilen ähnliche Abbildung der Ellipse auf den Einheitskreis, d. h. das so definierte $t$ stimmt mit der ebenso bezeichneten Größe der artt. [3.] und [4.] überein. Insbesondere ist die Reihe [14] mit [7] identisch, womit für die letztere Reihe das »Gesetz gefunden« ist. Der Übergang von den Logarithmen zu den Zahlen führt zu der Produktdarstellung [15] von $t$, die unmittelbar in die üblichen Bezeichnungen der Theorie der elliptischen Funktionen übertragen werden kann. Setzt man nämlich wie oben S. 321
\[\lambda^2=\mu^4=x=q,\quad\lambda u^2=\mu^2u^2=y=e^{2\pi iv},\]
\noindent so ist
\[t=\frac{P(x,y)}{Q(x,y)}=\frac{\vartheta_{00}(v\mid q)}{\vartheta_{10}(v\mid q)}\]
\noindent und der Übergang von [15] zu [16] vollzieht sich durch Anwendung der oft genannten Identität zwischen den Produkt- und Reihenformen der Thetafunktionen, siehe z. B. die Gleichung $\#$ S. 293.
\par\indent In der Gleichung [18], bezw. der berichtigten Form [18a], S. 319, und ebenso in [21] hat $T$ die Bedeutung
\[T=u+\lambda u^{-1},\]
\noindent während die in den artt. [3.] und [4.] mit $T$ bezeichnete Größe in [21] mit $T'$ bezeichnet ist. In [22] ist $T$ wieder in der Bedeutung $U+\lambda U^{-1}$ zu nehmen, wo also, vergl. art. [5.], S. 315,
\[U=\cos E+i\sin E\]
\noindent sich auf einen andern Ellipsenpunkt bezieht als $u$. Aus [22] ergibt sich für den reziproken Wert des in Gl. [24] mit $A$ bezeichneten Faktors zunächst die in [23] unter dem Logarithmuszeichen auftretende Produktdarstellung, die vermöge der schon in der Scheda Ac (siehe oben S. 201, art. [9.], Gl. 3), wo $z=\lambda^2$ zu nehmen ist) angegebenen Umformung in den reziproken Wert von
\[(1+\lambda^2+\lambda^6+\lambda^{12}+{.}{.}{.})^2\]
\noindent verwandelt werden kann. Damit ist nun auch für die in der Gl. [5a], oben S. 314 auftretende Reihe $\ast$ das »Gesetz gefunden«. Im art. [8.] ist in den sonst üblichen Bezeichnungen
\[\begin{aligned}
a&=p(\mu^4)=p(\lambda^2)=\vartheta_{00}(0\mid\lambda^2),\\
b&=r(\mu^4)=r(\lambda^2)=\vartheta_{10}(0\mid\lambda^2),
\end{aligned}\]
\noindent und nun folgt aus der bekannten Formel
\[2\vartheta_{00}(0\mid q^2)\vartheta_{10}(0\mid q^2)=\vartheta_{10}^2(0\mid q)\]
\noindent die im art. [6.], Gl. [20] angegebene Formel
\par\bigskip\noindent\rule{.25\linewidth}{.4pt}\par\smallskip
{\small\par\indent **) Wir machen besonders darauf aufmerksam, daß \textsc{Gauss} im Texte die Bezeichnung \textls[120]{Potential} benutzt und zwar in einer Weise, die darauf schließen läßt, daß ihm dieser Kunstausdruck damals (1839) völlig geläufig war. In dem älteren Teile der Abhandlung artt. [1.]---[4.] findet er sich nicht.\par}
\clearpage\SourceLeaf{332}\ArithmeticHead{325}{R}{{\scriptsize BEMERKUNGEN ZUR THEORIE DER TRANSZENDENTEN FUNKTIONEN.}}
\[ab=\tfrac12(r(\lambda))^2=\tfrac12\vartheta_{10}^2(0\mathord{\raisebox{-.3ex}{\rule{.3pt}{.6ex}}}\lambda).\]
\noindent Ferner folgt aus [18a], wenn man
\[\tau^2=\frac{a^2}{b^2}t^2,\quad k=\frac{b^2}{a^2}=\frac{\vartheta_{10}^2(0\ \lambda^2)}{\vartheta_{00}^2(0\mid\lambda^2)}\]
\noindent setzt,
\[\frac1{a^2}\int_0^\tau\frac{d\tau}{\sqrt{(1-\tau^2)(1-k^2\tau^2)}}=\int_0^T\frac{dT}{\sqrt{4\mu^2-T^2}}=\arcsin\frac T{2\mu}=i\log u,\]
\noindent also in den Bezeichnungen von \textsc{Jacobi}
\[t=\sqrt{k}\,\sin\operatorname{am}\left(\frac{2K}\pi\arcsin\frac T{2\mu}\right)=\sqrt{k}\,\sin\operatorname{am}\left(\frac{2K}\pi i\log u\right),\]
\noindent wo
\[a^2=\vartheta_{00}^2(0\mid\lambda^2)=\frac{2K}\pi=\frac2\pi\int_0^1\frac{d\tau}{\sqrt{1-\tau^2)(1-k^2\tau^2)}},\quad\lambda=\mu^2\]
\noindent ist. Diese Formel befindet sich in Übereinstimmung mit derjenigen Lösung der hier in Rede stehenden Aufgabe, die H. A. \textsc{Schwarz} in der Abhandlung \textit{Notizia sulla rappresentazione conforme di un’ area ellittica sopra un’ area circolare}, Annali di Matematica pura ed applicata, 2. serie, \textbf{3} (1870), S. 166, Gesammelte math. Abhandlungen II. 1890, S. 102, gegeben hat.
\begin{flushright}\textsc{Schlesinger.}\end{flushright}
\par\bigskip\begin{center}\rule{.25\linewidth}{.4pt}\end{center}
\clearpage
\endgroup
